% ==========================================================================
% OG-Core x CLEWS/OSeMOSYS — how the two models work together: narrative deck.
% Follows presentation/NARRATIVE.md. Reuses two standalone diagram PDFs
% (diagrams/architecture.pdf, diagrams/loop.pdf); the rest is native Beamer.
% Build: latexmk -pdf slides.tex   (or see ../build.sh)
% ==========================================================================
\documentclass[aspectratio=169,11pt]{beamer}

\usepackage[T1]{fontenc}
\usepackage{helvet}
\renewcommand{\familydefault}{\sfdefault}
\usepackage{microtype}
\usepackage{amsmath,amssymb}
\usepackage{graphicx}
\usefonttheme{professionalfonts}

% OG-Core x CLEWS — shared color palette.
% Single source of truth, mirrored from ogclews_link/style.py (the matplotlib figure
% theme) so the TikZ diagrams, the Beamer deck, and the data figures all share one
% editorial, colorblind-safe palette. \input this from any preamble; it only defines
% colors (no package loads), so it is safe in both standalone and beamer documents.

\definecolor{ink}{HTML}{222222}      % primary text / heavy rules
\definecolor{sub}{HTML}{555555}      % secondary text
\definecolor{mute}{HTML}{888888}     % tertiary / source lines
\definecolor{grid}{HTML}{E6E6E6}     % faint rules, fills

% the two models + policy — the load-bearing semantic colors
\definecolor{ogc}{HTML}{0F5499}      % OG-Core      (FT oxford blue)
\definecolor{clews}{HTML}{0D7680}    % CLEWS/OSeMOSYS (teal)
\definecolor{policy}{HTML}{E3120B}   % policy lever (Economist red)

% accents / signed values (match style.py)
\definecolor{claret}{HTML}{990F3D}   % kicker rule
\definecolor{gain}{HTML}{2166AC}     % positive change (blue)
\definecolor{loss}{HTML}{B2182B}     % negative change (red)
\definecolor{amber}{HTML}{E69F00}    % categorical 4
\definecolor{purple}{HTML}{7A6FAC}   % categorical 5
\definecolor{slate}{HTML}{7F8C8D}    % categorical 6 / neutral

% sequential blue ramp (ordered income groups, poorest -> richest)
\definecolor{seqA}{HTML}{C6DBEF}
\definecolor{seqB}{HTML}{9ECAE1}
\definecolor{seqC}{HTML}{6BAED6}
\definecolor{seqD}{HTML}{4292C6}
\definecolor{seqE}{HTML}{2171B5}
\definecolor{seqF}{HTML}{08519C}
\definecolor{seqG}{HTML}{08306B}


% ---- a minimal editorial theme keyed to the figure palette ------------------
\setbeamertemplate{navigation symbols}{}
\setbeamercolor{normal text}{fg=ink}
\setbeamercolor{frametitle}{fg=ink}
\setbeamercolor{structure}{fg=ogc}
\setbeamercolor{itemize item}{fg=ogc}
\setbeamercolor{itemize subitem}{fg=clews}
\setbeamercolor{enumerate item}{fg=ogc}
\setbeamercolor{block title}{fg=white,bg=ogc}
\setbeamercolor{block body}{fg=ink,bg=ogc!6}
\setbeamerfont{frametitle}{series=\bfseries,size=\large}
\setbeamerfont{title}{series=\bfseries}
\setbeamertemplate{itemize item}{\small\raise0.5pt\hbox{\textbullet}}
\setbeamertemplate{itemize subitem}{\footnotesize\raise0.5pt\hbox{--}}

% frametitle: a claret kicker rule above a bold ink title (mirrors style.py title_block)
\setbeamertemplate{frametitle}{%
  \vspace{0.42cm}%
  {\color{claret}\rule{0.62cm}{2.2pt}}\par\vspace{2.5pt}%
  \usebeamerfont{frametitle}\insertframetitle\par}

% spare footline: caption left, frame number right (full-width box so \hfill works)
\setbeamertemplate{footline}{%
  \hbox to\paperwidth{\hspace{0.5cm}%
    {\color{mute}\scriptsize OG-Core\,$\times$\,CLEWS · how the two models work together}%
    \hfill{\color{mute}\scriptsize\insertframenumber\,/\,\inserttotalframenumber}%
    \hspace{0.5cm}}%
  \vspace{3pt}}

% a titled diagram frame: title + diagram (sized to leave room) + a short 1-2 line note
\newcommand{\dgmframe}[3]{%
  \begin{frame}{#1}%
    \centering
    \includegraphics[width=\linewidth,height=0.66\textheight,keepaspectratio]{diagrams/#2.pdf}\par
    \vspace{4pt}{\scriptsize\color{mute}#3}%
  \end{frame}}

\title{How OG-Core and CLEWS work together}
\subtitle{What passes between them, and what it means}
\date{\today}

\begin{document}

% ---- title ------------------------------------------------------------------
\begin{frame}[plain]
  \vfill
  {\color{claret}\rule{1.2cm}{2.5pt}}\par\vspace{10pt}
  {\bfseries\LARGE\color{ink}How OG-Core and CLEWS work together}\par\vspace{8pt}
  {\large\color{sub}What passes between the economy and the energy system,\\ and what it means}\par\vspace{18pt}
  {\footnotesize\color{mute}OG-PHL calibration\, ·\, \texttt{ogclews-link}\, ·\, \today}\par
  \vfill
\end{frame}

% ---- the gap ----------------------------------------------------------------
\begin{frame}{Two analyses that usually don't meet}
  Energy-transition pathways and their economic consequences are normally studied with
  \textbf{separate models that don't share inputs}.
  \medskip
  \begin{itemize}\setlength\itemsep{7pt}
    \item A least-cost \textbf{energy} model gives build costs, electricity prices, and emissions —
          but not what they do to wages, the public budget, household welfare, or who gains and loses.
    \item A \textbf{macroeconomic} model gives those — but treats energy as a price and a quantity,
          with no representation of the system that supplies it.
  \end{itemize}
  \medskip
  {\color{ink}This couples the two}, so a change on the energy side carries through to the economy,
  and the economy's response carries back.
\end{frame}

% ---- the two models ---------------------------------------------------------
\dgmframe{What each model covers}{architecture}{%
  Each model is detailed exactly where the other is blank; coupling is the exchange of a few quantities
  between them.}

% ---- the linking price ------------------------------------------------------
\begin{frame}{The link is a price — the marginal cost of energy}
  The central quantity passed to the economy is the \textbf{energy price}, and it is not assumed or
  marked up.
  \medskip
  \begin{itemize}\setlength\itemsep{7pt}
    \item It is the \textbf{marginal cost of supply} in the least-cost solution — the cost of the
          \emph{next} unit of energy — which CLEWS produces directly (its shadow price), not as a markup.
    \item Marginal, not average, cost is the right signal: it is what an additional unit actually costs
          the system, so it is what should govern how much energy the economy chooses to use.
  \end{itemize}
  \medskip
  {\color{sub}That is what reaches households — not the average cost of everything produced.}
\end{frame}

% ---- the channels (overview) ------------------------------------------------
\begin{frame}{What connects them: eight channels}
  \textbf{\color{clews}From the energy system to the economy}\\[2pt]
  {\small the energy price\, ·\, the public grid build\, ·\, how capital-heavy the industry is\, ·\, air pollution and health}
  \medskip\par
  \textbf{\color{policy}Policy levers, applied to one or both models}\\[2pt]
  {\small a carbon price\, ·\, a tax credit for clean investment}
  \medskip\par
  \textbf{\color{ogc}From the economy back to the energy system}\\[2pt]
  {\small the energy demand to plan for\, ·\, the interest rate to discount with}
  \vfill
  {\footnotesize\color{sub}Each changes a \emph{different} part of the economy. They don't feed one
   another — their combined effect is worked out when the economy is solved.}
\end{frame}

% ---- how they reach a consistent answer -------------------------------------
\dgmframe{Reaching a consistent answer}{loop}{%
  Run the two in turn until the economy's energy demand and the system's supply agree — the consistent answer.}

% ---- channels act together --------------------------------------------------
\begin{frame}{The channels act together, not in sequence}
  Each channel sets a different parameter, and the economy is then solved \textbf{once} with all of them
  in place. Their interactions — cheaper energy changing labour supply, a capital subsidy moving the
  interest rate — are resolved \textbf{inside that single equilibrium}, not summed by hand. Because each
  touches a different parameter, the order they are applied in does not change the result.
  \vfill
  {\footnotesize\color{sub}Two pairs do act on the same quantity, so they are not combined:
   \textbf{energy price} and \textbf{carbon} both raise the price of energy; \textbf{capital intensity}
   and the \textbf{clean-investment tax credit} both move energy-sector capital, in opposite directions.}
\end{frame}

% ---- spotlight: capital intensity -------------------------------------------
\begin{frame}{A counterintuitive one: capital intensity}
  A clean fleet (wind, solar, carbon-capture) is \textbf{equipment-heavy} — more of its cost is the
  equipment built once, less is the fuel and labour paid for continuously, so a larger share of the
  industry's cost goes to capital.
  \medskip
  \begin{itemize}\setlength\itemsep{6pt}
    \item But energy is small and \textbf{price-inelastic}: making it cheaper to produce barely raises
          demand, so the industry doesn't expand. It meets the same demand more cheaply, and ends up
          using \textbf{\color{loss}\emph{less} total capital}, with the interest rate roughly flat.
  \end{itemize}
  \medskip
  {\color{sub}So this is a lever on the \textbf{price of energy} and the \textbf{capital--labour split} —
   not a way to draw capital in. (Drawing capital in is the tax credit, which pushes the other way.)}
\end{frame}

% ---- spotlight: health ------------------------------------------------------
\begin{frame}{Cleaner air: small on lives, real on productivity}
  The channel takes the change in fine-particle pollution (PM2.5), then scales it down sharply — the
  power sector is only about a tenth of the pollution people breathe, and the health response is
  concave. For the Philippines, roughly \textbf{8\%} of the raw change survives.
  \medskip
  \begin{itemize}\setlength\itemsep{6pt}
    \item \textbf{Mortality:} fewer early deaths — but they fall mostly on the elderly, so the effect on
          output is small.
    \item \textbf{Productivity:} working-age people are sick less and produce more — \textbf{this is where
          the output gain comes from}.
  \end{itemize}
\end{frame}

% ---- worked example ---------------------------------------------------------
\begin{frame}{One quantity around the loop}
  A scenario makes electricity cheaper to supply — say about \textbf{4\%} cheaper.
  \medskip
  \begin{enumerate}\setlength\itemsep{5pt}
    \item The lower cost flows through the \textbf{energy-price} channel and eases the price households pay.
    \item The economy is solved: households use more energy, and the relief is proportionally larger for
          \textbf{poorer households} — reported group by group.
    \item The \textbf{demand} channel passes that higher demand back to the energy system.
    \item The system re-solves at the higher demand; its mix and price shift; the economy responds again
          — until demand and supply line up.
  \end{enumerate}
  \medskip
  {\footnotesize\color{mute}One full circuit:\ \ \textbf{\color{clews}price} $\rightarrow$
   \textbf{\color{ogc}demand} $\rightarrow$ \textbf{\color{clews}price} $\rightarrow \cdots$}
\end{frame}

% ---- close ------------------------------------------------------------------
\begin{frame}{What the coupling adds}
  Neither model alone produces this:
  \medskip
  \begin{itemize}\setlength\itemsep{7pt}
    \item \textbf{Who} a change in the energy system helps or hurts — the effect by income group.
    \item \textbf{What cleaner air is worth} — in lives, and in a more productive workforce.
    \item A \textbf{mutually consistent} energy--economy outcome, rather than each side assuming the other.
  \end{itemize}
  \vfill
  {\footnotesize\color{mute}Plain-language guide: \texttt{presentation/NARRATIVE.md}\, ·\, OG-PHL
   calibration\, ·\, \texttt{ogclews-link}}
\end{frame}

\end{document}
